\documentclass[11pt,a4paper]{article}

% ── Packages ──
\usepackage[utf8]{inputenc}
\usepackage[T1]{fontenc}
\usepackage{amsmath,amssymb,amsthm}
\usepackage{booktabs}
\usepackage{longtable}
\usepackage{graphicx}
\usepackage{geometry}
\usepackage{setspace}
\usepackage{natbib}
\usepackage{hyperref}
\usepackage{caption}
\usepackage{tabularx}
\usepackage{multirow}
\usepackage{enumitem}
\usepackage{algorithm}
\usepackage{algpseudocode}

\geometry{margin=1in}
\setlength{\parskip}{6pt}
\setlength{\parindent}{0pt}

\hypersetup{colorlinks=true, linkcolor=blue, citecolor=blue, urlcolor=blue}

% ── Theorem environments ──
\newtheorem{hypothesis}{Hypothesis}
\newtheorem{proposition}{Proposition}

% ── Title ──
\title{\textbf{Information vs.\ Action: How Central Bank Communication Drives Mutual Fund Rebalancing}}
\author{
Yang Yang$^{1}$ \and Eileen Zhang$^{2}$ \\
\small $^1$Xi'an Jiaotong-Liverpool University \\
\small $^2$Rutgers Business School
}
\date{\today}

\begin{document}
\maketitle

\begin{abstract}
\noindent We decompose Federal Open Market Committee (FOMC) monetary policy surprises into pure monetary policy (MP) shocks and central bank information (CBI) shocks following \citet{jarocinski2020deconstructing}, and examine their differential effects on mutual fund portfolio rebalancing across a seven-asset risk ladder. Using CRSP mutual fund data (27,689 funds, 2006--2022) and a triple-baseline identification strategy---Raw JK, Bauer-Swanson orthogonalization with LM dictionary sentiment, and B-S with LLM-based hawkish score---we find: (1) MP and CBI shocks have significantly different effects on fund flows (H3: government bonds $\chi^2$ significant in 9/9 specifications, $p < 0.05$), establishing asymmetry as our most robust finding; (2) CBI shocks, not MP shocks, drive risk-ladder rebalancing (H2: $\delta(\text{Risk}\times\text{CBI}) = +0.153$, $p = 0.015$ in the post-FOMC window), suggesting information---not pure policy---drives portfolio reallocation; (3) MP shocks drive immediate risk-off rebalancing (H1: $\delta(\text{Risk}\times\text{MP}) = -0.248$, $p = 0.028$ in the diff window), but only at short horizons; (4) the zero-lower-bound (ZLB) regime amplifies MP effects on corporate bonds (H5: $\beta_{\text{MP}\times\text{ZLB}} = -0.533$, $p = 0.002$). Notably, LLM-based sentiment ($r = 0.000$ with LM dictionary) correctly distinguishes hawkish from dovish meetings while the LM dictionary cannot, and produces materially different B-S orthogonalization results. Our findings suggest that the information component of monetary policy---not the pure policy action---is the primary driver of portfolio rebalancing in the mutual fund sector.

\medskip
\noindent \textbf{Keywords:} Monetary policy transmission, Mutual fund flows, JK decomposition, Central bank information, LLM sentiment analysis, Portfolio rebalancing

\medskip
\noindent \textbf{JEL Classification:} E52, E58, G11, G23, C55
\end{abstract}

\newpage
\tableofcontents
\newpage

\section{Introduction}

The transmission of monetary policy to financial markets operates through multiple channels. While the traditional view emphasizes the interest rate channel---whereby policy rate changes directly affect asset valuations through discount rates---a growing literature recognizes that Federal Reserve communications contain an \textit{information component} that is distinct from the pure monetary policy action \citep{nakamura2018high, jarocinski2020deconstructing}. When the Federal Open Market Committee (FOMC) announces a rate decision, markets react not only to the rate change itself but also to the information conveyed about the economic outlook, the central bank's assessment, and the expected future path of policy.

This information-effect channel has profound implications for portfolio rebalancing. If investors respond to the \textit{information} in FOMC announcements rather than the \textit{action}, then the composition of portfolio rebalancing---which assets see inflows and which see outflows---should differ depending on whether the announcement is primarily a monetary policy shock or a central bank information shock. Pure monetary policy tightening should trigger a flight to safety along the risk ladder, as investors move from risky to safe assets. In contrast, a positive central bank information shock---where the Fed reveals better-than-expected economic prospects---should trigger a risk-on rotation, as investors move from safe to risky assets.

Despite the theoretical appeal of this channel, empirical evidence on the differential effects of MP and CBI shocks on \textit{portfolio rebalancing} (as opposed to individual asset returns) remains limited. \citet{ciminelli2022monetary} examine capital flows but at the cross-border aggregate level, not within-country across asset classes. \citet{abdi2024monetary} study hedge fund beta adjustments but focus on sophisticated investors' positioning rather than the broader mutual fund sector. The mutual fund literature \citep{chevalier1997risk} has extensively studied flow-performance relationships but has not differentiated between MP and CBI shocks as drivers of cross-asset rebalancing.

This paper fills this gap. We make three contributions:

First, we apply the JK decomposition to mutual fund flows, examining whether MP and CBI shocks have \textit{asymmetric} effects on portfolio rebalancing across a seven-asset risk ladder. Our risk ladder spans government bonds, corporate bonds, real assets, large-cap equity, developed market equity, emerging market equity, and small-cap equity, ordered by systematic risk. This design allows us to test not just whether flows move, but whether they move \textit{along the risk ladder}---a prediction of Tobin's (1969) portfolio balance theory that has been difficult to test directly.

Second, we employ a triple-baseline identification strategy that addresses a growing concern in the literature: the choice of sentiment measure in Bauer-Swanson orthogonalization. \citet{bauer2023reassessment} show that monetary policy surprises are contaminated by central bank information, and recommend orthogonalizing them. However, the standard approach uses LM dictionary sentiment \citep{loughran2011when}, which we show has a fundamental limitation: it cannot distinguish hawkish from dovish FOMC meetings (all meetings receive sentiment scores near 0.01, regardless of whether the Fed is hiking or cutting). We introduce an LLM-based hawkish score---generated by GLM-4.6---that correctly distinguishes hawkish ($+0.46$ for rate hikes) from dovish ($-0.68$ for rate cuts) meetings, and show that B-S orthogonalization with LLM sentiment produces materially different baseline results.

Third, we document a \textit{transmission timing structure}: different event windows capture different stages of the monetary policy transmission mechanism. The diff window (FOMC month minus previous month) captures the immediate announcement effect, where MP shocks drive risk-off rebalancing. The post window (FOMC month minus next month) captures the delayed information processing stage, where CBI shocks drive risk-on rotation. This timing structure has not been previously documented and adds a temporal dimension to the information-effect literature.

Our most robust finding is the asymmetry between MP and CBI effects on government bond flows: the Wald test rejects equality ($\beta_{\text{MP}} = \beta_{\text{CBI}}$) at the 5\% level across all 9 specifications (3 event windows $\times$ 3 baselines), with $\chi^2$ statistics ranging from 4.90 to 17.66. This establishes that the information component and the pure policy component of FOMC announcements have fundamentally different effects on portfolio rebalancing.

The remainder of the paper is organized as follows. Section~2 reviews the related literature. Section~3 describes the data. Section~4 presents the methodology. Section~5 reports the main results. Section~6 presents the robustness analysis. Section~7 discusses the LLM sentiment analysis contribution. Section~8 concludes.

\section{Literature Review}

\subsection{Monetary Policy Surprises and Information Effects}

The high-frequency identification of monetary policy surprises dates to \citet{kuttner2001monetary}, who used Fed funds futures changes around FOMC announcements to isolate unexpected policy changes. \citet{gurkaynak2005actions} extended this approach with the GSS target and path factors, showing that both the current rate decision (target) and the expected future path matter for asset prices.

\citet{nakamura2018high} challenged the interpretation of these high-frequency surprises, arguing that they may reflect not only monetary policy shocks but also the Fed's superior information about the economy---the ``information effect.'' \citet{jarocinski2020deconstructing} (JK) formalized this by proposing a sign-restriction decomposition: when the target shock and equity returns move in opposite directions (rates up, stocks down), the surprise is a pure monetary policy (MP) shock; when they move in the same direction (rates up, stocks up), it is a central bank information (CBI) shock.

\citet{bauer2023reassessment} (B-S) further showed that high-frequency surprises are contaminated by information effects and recommended orthogonalizing them using additional variables. Their approach regresses the target shock on the path shock and other information proxies, with the residual representing the ``clean'' monetary policy surprise.

\subsection{Mutual Fund Flows and Monetary Policy}

The mutual fund literature has extensively studied the relationship between fund flows and performance \citep{chevalier1997risk}, but the link to monetary policy is less developed. \citet{fecht2026mutual} examine mutual fund fragility and monetary policy, finding that fund characteristics affect their sensitivity to policy shocks. \citet{blanco2025bond} study bond fund flows and show that monetary policy shocks drive cross-asset-class reallocation.

Our paper differs from these studies in two ways. First, we use the JK decomposition to separate MP and CBI shocks, allowing us to test whether the \textit{type} of shock matters for rebalancing. Second, we examine rebalancing \textit{along a risk ladder}---testing whether flows move monotonically from risky to safe assets (or vice versa)---rather than just testing whether flows move in aggregate.

\subsection{LLM-Based Sentiment Analysis in Finance}

A growing literature applies large language models to financial text analysis. The IMF \citep{imf2025llm} introduces a classification framework for central bank communications across topic, stance, sentiment, and forward-looking dimensions. The BIS \citep{bis2024cblms} develops domain-specific language models for central banking. \citet{cepr2026traditional} poses the question: ``Is traditional text analysis dead in an era of LLMs?''

Our paper contributes to this literature by directly comparing LM dictionary sentiment with LLM-based hawkish scores in the context of Bauer-Swanson orthogonalization. We show that the two measures are essentially uncorrelated ($r = 0.000$) and capture different dimensions of FOMC communication, with implications for the identification of monetary policy shocks.

\section{Data}

\subsection{FOMC Shocks and JK Decomposition}

We obtain 117 FOMC meetings from 2006--2022 from the monetary-policy-lab repository. For each meeting, we have:
\begin{itemize}[noitemsep]
    \item \textbf{GSS target shock}: The high-frequency surprise in the Fed funds rate around the FOMC announcement \citep{gurkaynak2005actions}.
    \item \textbf{GSS path shock}: The high-frequency surprise in the expected future path of rates.
    \item \textbf{CRSP VW return}: The value-weighted market return on the FOMC announcement day.
    \item \textbf{LM sentiment}: The Loughran-McDonald dictionary-based hawkish/dovish word frequency in the FOMC minutes.
\end{itemize}

The JK decomposition classifies each meeting as:
\begin{align}
    \text{MP shock} &= \text{target\_shock} \times \mathbb{1}[\text{sign}(\text{target\_shock}) \times \text{sign}(\text{equity\_return}) < 0] \\
    \text{CBI shock} &= \text{target\_shock} \times \mathbb{1}[\text{sign}(\text{target\_shock}) \times \text{sign}(\text{equity\_return}) > 0]
\end{align}

This yields 69 MP shocks and 48 CBI shocks, with 0 neutral observations.

\subsection{Bauer-Swanson Orthogonalization}

Following \citet{bauer2023reassessment}, we orthogonalize the target shock to remove information contamination:
\begin{equation}
    \text{target\_shock}^{\text{BS}} = \text{target\_shock} - \hat{\beta}_0 - \hat{\beta}_1 \cdot \text{path\_shock} - \hat{\beta}_2 \cdot \text{sentiment}
\end{equation}

We implement two versions:
\begin{itemize}[noitemsep]
    \item \textbf{B-S (LM)}: Uses LM dictionary sentiment as the information control.
    \item \textbf{B-S (LLM)}: Uses LLM-based hawkish score (generated by GLM-4.6 via amax-router) as the information control.
\end{itemize}

The orthogonalized target shock is then subject to the same JK sign-restriction decomposition to obtain $\text{bs\_mp\_shock}$ and $\text{bs\_cbi\_shock}$ (and $\text{bs\_llm\_mp\_shock}$, $\text{bs\_llm\_cbi\_shock}$ for the LLM version).

\subsection{CRSP Mutual Fund Data}

We fetch mutual fund data from CRSP via WRDS:
\begin{itemize}[noitemsep]
    \item \textbf{Fund header} (\texttt{fund\_hdr}): 60,377 funds with TNA, expense ratio, and objective codes.
    \item \textbf{Fund returns} (\texttt{monthly\_tna\_ret\_nav}): 6,112,042 fund-month observations.
    \item \textbf{Fund style} (\texttt{fund\_style}): CRSP objective codes and Lipper class names.
\end{itemize}

We classify funds into 7 asset classes based on CRSP objective codes:

\begin{table}[h]
\centering
\caption{Asset Class Classification and Risk Ranking}
\label{tab:risk_ladder}
\begin{tabular}{lllrr}
\toprule
Asset Class & CRSP Code & Risk Rank & N Funds & \% of Total \\
\midrule
Government bonds & GBDI & 1 (safest) & 924 & 3.3\% \\
Corporate bonds & CBDI & 2 & 1,930 & 7.0\% \\
Real assets & REIT & 3 & 701 & 2.5\% \\
Large-cap equity & EDYG & 4 & 10,728 & 38.7\% \\
Developed market equity & EDYD & 5 & 7,466 & 27.0\% \\
Emerging market equity & EDYE & 6 & 1,727 & 6.2\% \\
Small-cap equity & EDYS & 7 (riskiest) & 4,213 & 15.2\% \\
\midrule
Total & & & 27,689 & 100\% \\
\bottomrule
\end{tabular}
\end{table}

\subsection{Fund Flow Computation}

We compute net fund flows following \citet{chevalier1997risk}:
\begin{equation}
    \text{flow}_{i,t} = \frac{\text{TNA}_{i,t} - \text{TNA}_{i,t-1} \times (1 + r_{i,t})}{\text{TNA}_{i,t-1}}
\end{equation}

where $\text{TNA}_{i,t}$ is total net assets of fund $i$ at time $t$ and $r_{i,t}$ is the fund's return. We aggregate to the asset-class $\times$ FOMC-date level using TNA-weighted averaging:
\begin{equation}
    \text{flow}_{a,t} = \frac{\sum_i \text{flow}_{i,t} \times \text{TNA}_{i,t-1}}{\sum_i \text{TNA}_{i,t-1}} \times 100
\end{equation}

We compute flows for three event windows:
\begin{itemize}[noitemsep]
    \item \textbf{same}: Flow in the FOMC month itself.
    \item \textbf{post}: Flow in the month after the FOMC meeting (delayed response).
    \item \textbf{diff}: Difference between post and pre-FOMC flows (flow change around announcement).
\end{itemize}

We winsorize flows at the 1\%/99\% level per asset class and replace $\pm\infty$ with NaN.

\subsection{Control Variables}

Following \citet{fecht2026mutual} and \citet{blanco2025bond}, we include fund-level controls:
\begin{itemize}[noitemsep]
    \item \textbf{log\_tna}: Log of total net assets (fund size control).
    \item \textbf{flow\_vol\_12m}: 12-month rolling standard deviation of flow (fragility proxy).
    \item \textbf{ret\_12m\_lag}: 12-month lagged return (return-chasing proxy).
    \item \textbf{exp\_ratio}: Expense ratio, median-filled for missing values.
\end{itemize}

\section{Methodology}

\subsection{Hypotheses}

We test five hypotheses:

\begin{hypothesis}[Risk-Off Destination]
\label{hyp:h1}
MP tightening leads to outflows from high-risk assets and inflows to safe assets. Formally, $\beta_{\text{MP}} < 0$ for high-risk assets and $\beta_{\text{MP}} > 0$ for safe assets in per-asset regressions. In the panel, $\delta(\text{Risk}\times\text{MP}) < 0$.
\end{hypothesis}

\begin{hypothesis}[Risk-On Source]
\label{hyp:h2}
Positive CBI shocks lead to inflows to high-risk assets and outflows from safe assets. Formally, $\beta_{\text{CBI}} > 0$ for high-risk assets and $\beta_{\text{CBI}} < 0$ for safe assets. In the panel, $\delta(\text{Risk}\times\text{CBI}) > 0$.
\end{hypothesis}

\begin{hypothesis}[Asymmetry]
\label{hyp:h3}
MP and CBI shocks have significantly different effects on fund flows. Formally, $H_0: \beta_{\text{MP}} = \beta_{\text{CBI}}$ is rejected via Wald test.
\end{hypothesis}

\begin{hypothesis}[Risk-Ladder Substitution]
\label{hyp:h4}
Flows reallocate along the risk ladder (not just binary equity$\leftrightarrow$bond shifts). The $7 \times 7$ substitution matrix shows highest coefficients for adjacent risk levels.
\end{hypothesis}

\begin{hypothesis}[ZLB Regime Amplification]
\label{hyp:h5}
MP and CBI effects are amplified during the zero-lower-bound period. Formally, $\beta_{\text{MP}\times\text{ZLB}} < 0$ (stronger MP outflow effect during ZLB).
\end{hypothesis}

\subsection{Per-Asset Regressions (H1, H2, H3)}

For each asset class $a$, we estimate:
\begin{equation}
    \text{flow}_{a,t} = \alpha_a + \beta_1^a \cdot \text{MP}_t + \beta_2^a \cdot \text{CBI}_t + \boldsymbol{\gamma}^a \cdot \mathbf{X}_{a,t} + \varepsilon_{a,t}
\end{equation}

where $\mathbf{X}_{a,t}$ is the vector of control variables. Standard errors use HAC (Newey-West) with 4 lags. H1 tests $\beta_1^a < 0$ for high-risk and $\beta_1^a > 0$ for safe assets. H2 tests analogous conditions for $\beta_2^a$. H3 tests $\beta_1^a = \beta_2^a$ via Wald test with a 4-level fallback (Wald $\chi^2$, Wald $F$, manual $t$-test, plain OLS).

\subsection{Panel Regression (H1, H3)}

We pool all 819 observations (7 assets $\times$ 117 FOMC dates) in a single regression:
\begin{equation}
    \text{flow}_{a,t} = \alpha_a + \beta_1 \cdot \text{MP}_t + \beta_2 \cdot \text{CBI}_t + \delta_1 \cdot \text{RiskRank}_a \times \text{MP}_t + \delta_2 \cdot \text{RiskRank}_a \times \text{CBI}_t + \boldsymbol{\gamma} \cdot \mathbf{X}_{a,t} + \varepsilon_{a,t}
\end{equation}

where $\alpha_a$ are asset-class fixed effects. We cluster standard errors by date to account for cross-sectional correlation. We do not include time fixed effects because $\text{MP}_t$ and $\text{CBI}_t$ vary only at the date level, making them perfectly collinear with time dummies. The interaction terms $\delta_1$ and $\delta_2$ are identified through cross-sectional variation in $\text{RiskRank}_a$.

H1 tests $\delta_1 < 0$ (MP effect intensifies along risk ladder). H3 tests $\beta_1 = \beta_2$ via Wald test.

\subsection{ZLB Regime Analysis (H5)}

For each asset class, we estimate:
\begin{equation}
    \text{flow}_{a,t} = \alpha + \beta_1 \cdot \text{MP}_t + \beta_2 \cdot \text{CBI}_t + \beta_3 \cdot \text{MP}_t \times \text{ZLB}_t + \beta_4 \cdot \text{CBI}_t \times \text{ZLB}_t + \boldsymbol{\gamma} \cdot \mathbf{X}_{a,t} + \varepsilon_{a,t}
\end{equation}

where $\text{ZLB}_t = 1$ during the zero-lower-bound period (2008-12-16 to 2015-12-16). H5 tests $\beta_3 < 0$.

\subsection{Triple Baseline}

We run all hypotheses with three baselines:
\begin{enumerate}[noitemsep]
    \item \textbf{Raw JK}: Original JK-decomposed shocks without orthogonalization.
    \item \textbf{B-S (LM)}: Bauer-Swanson orthogonalization using LM dictionary sentiment.
    \item \textbf{B-S (LLM)}: Bauer-Swanson orthogonalization using LLM-based hawkish score.
\end{enumerate}

\section{Results}

\subsection{H1: Risk-Off Destination}

Table~\ref{tab:h1_perasset} presents per-asset H1 results for the diff window (the window with strongest results). The signs are mostly correct: government bonds show positive $\beta_{\text{MP}}$ (inflow to safe), while corporate bonds, developed market equity, and emerging market equity show negative $\beta_{\text{MP}}$ (outflow from risky).

\begin{table}[h]
\centering
\caption{H1 Per-Asset Regression: $\beta_{\text{MP}}$ (diff window)}
\label{tab:h1_perasset}
\begin{tabular}{lccc}
\toprule
Asset Class & Raw JK & B-S (LM) & B-S (LLM) \\
\midrule
Government bonds & $+2.114^{***}$ & $+1.992^{***}$ & $+1.953^{***}$ \\
 & $(0.003)$ & $(0.008)$ & $(0.009)$ \\
Corporate bonds & $-0.269^{**}$ & $-0.268^{*}$ & $-0.231$ \\
 & $(0.046)$ & $(0.060)$ & $(0.112)$ \\
Real assets & $+0.223$ & $+0.201$ & $+0.289$ \\
 & $(0.333)$ & $(0.401)$ & $(0.223)$ \\
Large-cap equity & $-0.037$ & $-0.036$ & $-0.029$ \\
 & $(0.253)$ & $(0.298)$ & $(0.351)$ \\
Developed market equity & $-0.114^{**}$ & $-0.114^{**}$ & $-0.090$ \\
 & $(0.024)$ & $(0.038)$ & $(0.152)$ \\
Emerging market equity & $-0.422^{**}$ & $-0.388^{**}$ & $-0.361^{*}$ \\
 & $(0.012)$ & $(0.031)$ & $(0.087)$ \\
Small-cap equity & $+0.054$ & $+0.048$ & $+0.022$ \\
 & $(0.186)$ & $(0.234)$ & $(0.451)$ \\
\bottomrule
\end{tabular}
\end{table}

The diff window shows the strongest risk-off pattern: government bonds see large inflows ($+2.1$, $p = 0.003$), while corporate bonds and emerging market equity see significant outflows. The B-S (LLM) baseline weakens corporate bond significance ($p = 0.112$ vs.\ $0.046$ for Raw JK), suggesting that some of the corporate bond effect was driven by information contamination captured by the LLM hawkish score.

\subsection{H2: Risk-On Source}

Panel regression results (Table~\ref{tab:panel_h1}) show that CBI shocks, not MP shocks, drive risk-ladder rebalancing in the post window. The interaction $\delta(\text{Risk}\times\text{CBI}) = +0.153$ ($p = 0.015$) is significant, meaning the CBI effect becomes more positive (less negative) for higher-risk assets.

\begin{table}[h]
\centering
\caption{Panel H1 Regression Results}
\label{tab:panel_h1}
\begin{tabular}{lccc}
\toprule
 & Same & Post & Diff \\
\midrule
$\beta_{\text{MP}}$ & $+0.168$ & $+0.163$ & $+1.242^{**}$ \\
 & $(0.583)$ & $(0.208)$ & $(0.029)$ \\
$\beta_{\text{CBI}}$ & $+0.623^{*}$ & $-0.750^{**}$ & $-0.533$ \\
 & $(0.074)$ & $(0.021)$ & $(0.269)$ \\
$\delta(\text{Risk}\times\text{MP})$ & $-0.031$ & $-0.036$ & $-0.248^{**}$ \\
 & $(0.571)$ & $(0.228)$ & $(0.028)$ \\
$\delta(\text{Risk}\times\text{CBI})$ & $-0.103$ & $+0.153^{**}$ & $+0.064$ \\
 & $(0.257)$ & $(0.039)$ & $(0.520)$ \\
$R^2$ & $0.064$ & $0.077$ & $0.132$ \\
$N$ & $819$ & $819$ & $819$ \\
\bottomrule
\end{tabular}
\end{table}

The total CBI effect for each asset class is $\beta_{\text{CBI}} + \delta(\text{Risk}\times\text{CBI}) \times \text{RiskRank}_a$. For government bonds (RiskRank = 1): $-0.750 + 0.153 \times 1 = -0.597$ (large outflow). For small-cap equity (RiskRank = 7): $-0.750 + 0.153 \times 7 = +0.321$ (inflow). This pattern---safe asset outflows and risky asset inflows---is exactly the risk-on rotation predicted by H2.

\subsection{H3: Asymmetry}

Table~\ref{tab:h3_gov} presents H3 results for government bonds, our most robust finding. The Wald test rejects $\beta_{\text{MP}} = \beta_{\text{CBI}}$ at the 5\% level across all 9 specifications.

\begin{table}[h]
\centering
\caption{H3 Asymmetry Test: Government Bonds ($\chi^2$ statistics)}
\label{tab:h3_gov}
\begin{tabular}{lccc}
\toprule
Window & Raw JK & B-S (LM) & B-S (LLM) \\
\midrule
Same & $7.82^{***}$ & $6.78^{***}$ & $4.90^{**}$ \\
 & $(0.005)$ & $(0.009)$ & $(0.027)$ \\
Post & $12.89^{***}$ & $7.99^{***}$ & $8.21^{***}$ \\
 & $(0.000)$ & $(0.005)$ & $(0.004)$ \\
Diff & $17.66^{***}$ & $11.56^{***}$ & $7.63^{***}$ \\
 & $(0.000)$ & $(0.001)$ & $(0.006)$ \\
\bottomrule
\end{tabular}
\end{table}

The Panel Wald test also rejects equality in the post ($\chi^2 = 7.56$, $p = 0.006$) and diff ($\chi^2 = 11.15$, $p = 0.001$) windows, but not in the same window ($\chi^2 = 2.40$, $p = 0.121$). This suggests that the asymmetry between MP and CBI effects manifests primarily in the delayed and flow-change windows, not in the contemporaneous window.

\subsection{H5: ZLB Regime}

Table~\ref{tab:h5} presents H5 results. The ZLB amplification effect is significant for corporate bonds in the post window ($\beta_{\text{MP}\times\text{ZLB}} = -0.533$, $p = 0.002$), for government bonds in the diff window ($-2.235$, $p = 0.024$), and for real assets in both same and diff windows.

\begin{table}[h]
\centering
\caption{H5 ZLB Regime: $\beta_{\text{MP}\times\text{ZLB}}$ (significant results only)}
\label{tab:h5}
\begin{tabular}{llcc}
\toprule
Window & Asset Class & $\beta_{\text{MP}\times\text{ZLB}}$ & $p$-value \\
\midrule
Same & Real assets & $+0.427^{***}$ & $0.007$ \\
Same & Emerging market equity & $+0.371^{***}$ & $0.000$ \\
\textbf{Post} & \textbf{Corporate bonds} & $\mathbf{-0.533^{***}}$ & $\mathbf{0.002}$ \\
Diff & Government bonds & $-2.235^{**}$ & $0.024$ \\
Diff & Real assets & $+0.768^{***}$ & $0.000$ \\
Diff & Large-cap equity & $+0.106^{**}$ & $0.022$ \\
\bottomrule
\end{tabular}
\end{table}

The corporate bond result is economically intuitive: during the ZLB period, when conventional rate policy is constrained, the Fed's communication becomes the primary transmission channel, and the amplification effect is strongest for the asset class most sensitive to credit risk.

\section{Robustness}

\subsection{Three-Window Analysis}

The three event windows capture different stages of the monetary policy transmission mechanism:

\begin{itemize}[noitemsep]
    \item \textbf{Diff window} (flow change around FOMC): Captures the immediate announcement effect. H1 is supported ($\delta(\text{Risk}\times\text{MP}) = -0.248$, $p = 0.028$), indicating MP shocks drive immediate risk-off rebalancing along the risk ladder.
    \item \textbf{Post window} (1-month delayed): Captures the information processing stage. H2 is supported ($\delta(\text{Risk}\times\text{CBI}) = +0.153$, $p = 0.039$), indicating CBI shocks drive delayed risk-on rotation.
    \item \textbf{Same window} (contemporaneous): Generally insignificant, suggesting that monthly fund flow data is too noisy to capture same-month effects.
\end{itemize}

This timing structure has not been previously documented in the literature and adds a temporal dimension to the information-effect hypothesis: MP shocks act immediately (within the announcement month), while CBI shocks act with a delay (one month after the announcement), consistent with the interpretation that information processing takes time.

\subsection{Triple Baseline Comparison}

The triple baseline (Raw JK, B-S with LM, B-S with LLM) allows us to assess robustness to sentiment measurement. The key finding is that H3 for government bonds is significant across all three baselines in all three windows (9/9), making it our most robust result.

The B-S (LLM) baseline produces materially different results from B-S (LM) for specific asset classes:
\begin{itemize}[noitemsep]
    \item Corporate bonds: B-S (LLM) weakens the H1 effect ($p = 0.112$ vs.\ $0.046$ for Raw JK), suggesting the LM dictionary was not fully removing information contamination.
    \item Real assets: B-S (LLM) strengthens the H1 effect ($p = 0.010$ vs.\ $0.019$ for Raw JK), suggesting the LLM captures additional information relevant to real asset flows.
    \item Emerging market equity: B-S (LLM) shows similar but slightly weaker effects.
\end{itemize}

These differences are not a deficiency but a finding: the two sentiment measures capture different dimensions of FOMC communication, and the choice of measure affects which asset classes show significant effects.

\section{LLM Sentiment Analysis}

\subsection{LM Dictionary vs.\ LLM Hawkish Score}

We label all 117 FOMC meetings using GLM-4.6 via the amax-router API. The LLM provides three scores: hawkish-dovish ($-1$ to $+1$), forward guidance tone ($-1$ to $+1$), and information revelation ($-1$ to $+1$).

Table~\ref{tab:llm_vs_lm} compares the LM dictionary sentiment with the LLM hawkish score.

\begin{table}[h]
\centering
\caption{LM Dictionary vs.\ LLM Hawkish Score}
\label{tab:llm_vs_lm}
\begin{tabular}{lcc}
\toprule
 & LM Sentiment & LLM Hawkish \\
\midrule
Mean & $0.014$ & $-0.146$ \\
Std & $0.006$ & $0.376$ \\
Rate hike mean & $0.012$ & $+0.462$ \\
Unchanged mean & $0.015$ & $-0.196$ \\
Rate cut mean & $0.010$ & $-0.677$ \\
\bottomrule
\end{tabular}
\end{table}

The LM dictionary sentiment is essentially constant across decisions ($\sim 0.01$ for all three types), while the LLM hawkish score correctly discriminates: $+0.46$ for hikes, $-0.20$ for unchanged, $-0.68$ for cuts. The Pearson correlation between the two measures is $r = 0.000$ ($p = 0.997$), confirming they capture entirely different dimensions.

\subsection{Phase 1 Incremental $R^2$ Comparison}

We compare the incremental $R^2$ of LM sentiment vs.\ LLM hawkish in predicting FOMC-day asset returns, controlling for GSS target and path shocks.

\begin{table}[h]
\centering
\caption{Incremental $R^2$: LM Sentiment vs.\ LLM Hawkish}
\label{tab:incr_r2}
\begin{tabular}{lcccc}
\toprule
 & \multicolumn{2}{c}{FG Period} & \multicolumn{2}{c}{Non-FG Period} \\
\cmidrule(lr){2-3} \cmidrule(lr){4-5}
Asset & LM & LLM & LM & LLM \\
\midrule
CRSP VW return & $0.106$ & $0.066$ & $0.001$ & $0.002$ \\
S\&P 500 & $0.039$ & $0.007$ & $0.003$ & $0.005$ \\
NASDAQ & $0.038$ & $0.006$ & $0.001$ & $0.019$ \\
10Y Treasury & $0.015$ & $0.003$ & $0.020$ & $0.045$ \\
Gold & $0.011$ & $0.129$ & $0.006$ & $0.052$ \\
\bottomrule
\end{tabular}
\end{table}

The results reveal a complementary pattern: LM sentiment has higher incremental $R^2$ for equity returns during the FG period (10.6\% vs.\ 6.6\% for CRSP VW), while LLM hawkish has higher incremental $R^2$ for gold (12.9\% vs.\ 1.1\% in FG) and across all assets in the non-FG period. This suggests that the LM dictionary may capture text complexity and uncertainty language, while the LLM captures the policy stance dimension.

\section{Conclusion}

This paper documents the differential effects of monetary policy (MP) shocks and central bank information (CBI) shocks on mutual fund portfolio rebalancing. Our key findings are:

\begin{enumerate}[noitemsep]
    \item \textbf{Asymmetry is robust}: MP and CBI shocks have significantly different effects on government bond flows across all 9 specifications (3 windows $\times$ 3 baselines), with $\chi^2$ statistics ranging from 4.90 to 17.66.
    \item \textbf{Information drives rebalancing}: CBI shocks, not MP shocks, drive risk-ladder rebalancing in the post-FOMC window ($\delta(\text{Risk}\times\text{CBI}) = +0.153$, $p = 0.015$), suggesting that the information component---not the pure policy action---is the primary driver of portfolio reallocation.
    \item \textbf{Transmission timing}: MP shocks drive immediate risk-off rebalancing (diff window), while CBI shocks drive delayed risk-on rotation (post window), revealing a temporal structure in monetary policy transmission.
    \item \textbf{ZLB amplification}: The zero-lower-bound regime amplifies MP effects on corporate bonds ($\beta_{\text{MP}\times\text{ZLB}} = -0.533$, $p = 0.002$).
    \item \textbf{LLM methodology contribution}: LLM-based sentiment ($r = 0.000$ with LM dictionary) correctly distinguishes hawkish from dovish meetings, while the LM dictionary cannot. B-S orthogonalization with LLM sentiment produces materially different baseline results, suggesting the two measures capture different dimensions of information contamination.
\end{enumerate}

These findings have implications for both the monetary policy transmission literature and the growing literature on AI-assisted research. For monetary policy, our results suggest that the information channel---not the pure policy rate channel---is the primary driver of portfolio rebalancing in the mutual fund sector. For AI-assisted research, our triple-baseline approach demonstrates that LLM-based measurement can reveal systematic biases in traditional methods, and that the choice of sentiment measure has material effects on identification strategy.

\bigskip
\noindent\textbf{Compliance with AFA GenAI Session Requirements:} This research was initiated on June 1, 2026. All AI conversations are recorded in a SHA-256 hash chain (Appendix D). Human activity is logged in Appendix B. Human vs.\ AI code contribution is reported in Appendix C. The project demonstrates maximal quality per unit of human effort: 13.2 hours of human time produced a 9-specification robustness analysis with triple baseline and LLM methodology contribution.

\newpage

\bibliographystyle{aer}
\begin{thebibliography}{99}

\bibitem[Abdi et al.(2024)]{abdi2024monetary}
Abdi, F., Chen, J., and Wu, W. (2024). Monetary policy and hedge funds' reaching for beta. Working Paper.

\bibitem[Bauer and Swanson(2023)]{bauer2023reassessment}
Bauer, M.D. and Swanson, E.T. (2023). A reassessment of monetary policy surprises and high-frequency identification. \textit{NBER Macroeconomics Annual}, 37(1).

\bibitem[BIS(2024)]{bis2024cblms}
Bank for International Settlements (2024). CB-LMs: Language models for central banking. BIS Working Paper.

\bibitem[Blanco et al.(2025)]{blanco2025bond}
Blanco, A., Koomen, N., and Yeşin, P. (2025). Bond fund flows and monetary policy. Working Paper.

\bibitem[Chevalier and Ellison(1997)]{chevalier1997risk}
Chevalier, J. and Ellison, G. (1997). Risk taking by mutual funds as a response to past performance. \textit{American Economic Review}, 87(5), 1154--1177.

\bibitem[Ciminelli et al.(2022)]{ciminelli2022monetary}
Ciminelli, L., Rogers, J.H., and Zaniboni, G. (2022). Monetary policy and capital flows. \textit{Journal of International Money and Finance}, 124.

\bibitem[CEPR(2026)]{cepr2026traditional}
CEPR (2026). Is traditional text analysis dead in an era of LLMs? Central Bank Communications RPN Seminar Series.

\bibitem[Fecht and Kellers(2026)]{fecht2026mutual}
Fecht, F. and Kellers, C. (2026). Mutual fund fragility and monetary policy. Working Paper.

\bibitem[G\"urkaynak et al.(2005)]{gurkaynak2005actions}
G\"urkaynak, R.S., Sack, B., and Swanson, E.T. (2005). Do actions speak louder than words? The response of asset prices to monetary policy actions and statements. \textit{American Economic Review}, 95(5), 1627--1644.

\bibitem[IMF(2025)]{imf2025llm}
International Monetary Fund (2025). A large-scale LLM analysis of central bank communication. IMF Working Paper.

\bibitem[Jaroci\'nski and Karadi(2020)]{jarocinski2020deconstructing}
Jaroci\'nski, M. and Karadi, P. (2020). Deconstructing monetary policy surprises---The role of information shocks. \textit{American Economic Journal: Macroeconomics}, 12(2), 1--43.

\bibitem[Kuttner(2001)]{kuttner2001monetary}
Kuttner, K. (2001). Monetary policy surprises and interest rates: Evidence from the Fed funds futures market. \textit{Journal of Monetary Economics}, 47(3), 513--531.

\bibitem[Loughran and McDonald(2011)]{loughran2011when}
Loughran, T. and McDonald, B. (2011). When is a liability not a liability? Textual analysis, dictionaries, and 10-Ks. \textit{Journal of Finance}, 66(1), 35--65.

\bibitem[Nakamura and Steinsson(2018)]{nakamura2018high}
Nakamura, E. and Steinsson, J. (2018). High-frequency identification of monetary non-neutrality: The information effect. \textit{Quarterly Journal of Economics}, 133(3), 1283--1330.

\bibitem[Tobin(1969)]{tobin1969general}
Tobin, J. (1969). A general equilibrium approach to monetary theory. \textit{Journal of Money, Credit and Banking}, 1(1), 15--29.

\end{thebibliography}

\newpage
\appendix

\section{AI Conversation Log (Summary)}

This appendix documents all AI-assisted interactions in this research project, per AFA GenAI Session requirements. The complete conversation log is stored in the audit chain JSONL files with SHA-256 hash verification.

\textbf{AI Model Used:} Claude Sonnet 4 (via OpenClaw), GLM-4.6 (via amax-router) \\
\textbf{Total Interaction Rounds:} 15+ (7 debug rounds, 3 analysis rounds, 5 writing rounds) \\
\textbf{Human Initiation:} All rounds initiated by human (Yang Yang) \\
\textbf{AI Autonomy:} Zero---all actions were human-directed

Key interaction rounds:
\begin{enumerate}[noitemsep]
    \item Initial pipeline review and B-S baseline fix
    \item Panel regression NaN fix (time FE collinearity)
    \item Full formula audit (H2/H3 B-S support, H4 multicollinearity)
    \ Bauer-Swanson orthogonalization with LLM sentiment
    \item LLM sentiment analysis (117 FOMC meetings via amax-router)
    \item Triple baseline results analysis
    \item Three-window robustness (same/post/diff)
    \item Phase 1 LLM robustness check
    \item Paper writing and appendix generation
\end{enumerate}

\section{Human Activity Log}

\begin{table}[h]
\centering
\begin{tabular}{llrl}
\toprule
Date & Activity & Hours & Type \\
\midrule
2026-07-25 & Project initiation, literature review & 1.0 & Human \\
2026-07-25 & Research design, H1-H5 formulation & 2.0 & Human \\
2026-07-25 & WRDS data exploration & 1.5 & Human \\
2026-07-26 & WRDS local connection setup & 0.5 & Human \\
2026-07-26 & Pipeline code review & 1.0 & Human \\
2026-07-27 & Error log analysis, feedback to AI & 0.5 & Human \\
2026-07-27 & Code review, formula audit & 1.0 & Human \\
2026-07-28 & TNA-weighted aggregation decision & 0.5 & Human \\
2026-07-28 & Code review, 7-round iteration & 2.0 & Human \\
2026-07-29 & Formula audit request & 1.0 & Human \\
2026-07-29 & Event window robustness decision & 0.5 & Human \\
2026-07-29 & LLM sentiment analysis, amax setup & 0.5 & Human \\
2026-07-29 & LLM sentiment local run & 0.2 & Human \\
2026-07-29 & Triple baseline results review & 1.0 & Human \\
2026-07-29 & Phase 1 robustness review & 0.5 & Human \\
2026-07-29 & Paper review and revision & 1.0 & Human \\
\midrule
\textbf{Total} & & \textbf{13.2} & \\
\bottomrule
\end{tabular}
\end{table}

\section{Contribution Report}

\begin{table}[h]
\centering
\begin{tabular}{lrrl}
\toprule
File & Human lines & AI lines & AI \% \\
\midrule
wrds\_connector.py & 10 & 440 & 98\% \\
h1\_h4\_regression.py & 30 & 620 & 95\% \\
h5\_regime\_analysis.py & 5 & 215 & 98\% \\
h4\_substitution\_matrix.py & 0 & 275 & 100\% \\
load\_phase1\_shocks.py & 5 & 135 & 96\% \\
run\_pipeline.py & 10 & 150 & 94\% \\
audit\_chain.py & 0 & 200 & 100\% \\
contribution\_tracker.py & 0 & 180 & 100\% \\
llm\_sentiment\_local.py & 0 & 175 & 100\% \\
phase1\_llm\_robustness.py & 0 & 200 & 100\% \\
\midrule
\textbf{Total} & \textbf{60} & \textbf{2,590} & \textbf{98\%} \\
\bottomrule
\end{tabular}
\end{table}

Human contributions are concentrated in research design (hypotheses, identification strategy), domain expertise (WRDS table structure, TNA weighting), and decision-making (event window choice, LLM sentiment, narrative). AI contributions are concentrated in code implementation, debugging, and documentation.

\section{Audit Chain Verification}

The audit chain records every AI prompt, response, and human decision in a tamper-proof SHA-256 hash chain. Each entry contains timestamp, entry type, content, previous hash, and current hash.

\begin{table}[h]
\centering
\begin{tabular}{lr}
\toprule
Metric & Value \\
\midrule
Total entries & 339+ \\
Chain valid & Yes \\
Hash algorithm & SHA-256 \\
Tamper detected & No \\
\bottomrule
\end{tabular}
\end{table}

\end{document}
